\SourceLeaf{61}
$\gamma$) His praecognitis facile coefficientium reliquorum expressiones finitae per scalam recursionis ante erutam deducuntur. Erat illa
\[
A^{(h+2)} = \frac{n(\nu-h)A^{(h)} - 2(h+1)A^{(h+1)}}{n(\nu+h+2)}
\]
adeoque h.\,l.\ ob $\nu=-1$
\[
A^{(h+2)} = \frac{-n(h+1)A^{(h)} - 2(h+1)A^{(h+1)}}{n(h+1)}
\]
sive
\[
A^{(h+2)} = -\Bigl(A^{(h)} + \frac{2}{n}A^{(h+1)}\Bigr).
\]
Si per hanc scalam calculos instituas, occurret lex simplicissima\HistoricalNote{*)}{*)}{Falsa est, vitio forsan typothetae, legis huius apud \textsc{Eulerum} l. c. \S\,275 expressio.}, fore
\[
A^{(h)} = (-1)^h\,\frac{2}{\BareRoot{(1-nn)}}\left(\frac{1-\BareRoot{(1-nn)}}{n}\right)^{h}.
\]
Ut illa demonstratione firmetur, sumamus valere ipsam pro coefficientibus usque ad $h^{\text{tum}}$, eritque per scalam recursionis
\[
\begin{aligned}
A^{(h+1)} & = (-1)^h\,\frac{2}{\BareRoot{(1-nn)}}\left(\frac{1-\BareRoot{(1-nn)}}{n}\right)^{h-1}\\&\quad + (-1)^{h+1}\,\frac{2}{n}\,\frac{2}{\BareRoot{(1-nn)}}\left(\frac{1-\BareRoot{(1-nn)}}{n}\right)^{h}
\end{aligned}
\]
sive, separato factore communi
\[
(-1)^h\,\frac{2}{\BareRoot{(1-nn)}}\left(\frac{1-\BareRoot{(1-nn)}}{n}\right)^{h-1}\left\{1-\frac{2}{n}\left(\frac{1-\BareRoot{(1-nn)}}{n}\right)\right\}.
\]
Est factor, uncinis inclusus
\[
\begin{aligned}
&= \frac{nn-2+2\BareRoot{(1-nn)}}{nn}\\&= -\frac{(1-nn)-2\BareRoot{(1-nn)}+1}{nn}\\&= -\left(\frac{1-\BareRoot{(1-nn)}}{n}\right)^{2},
\end{aligned}
\]
quo adiecto habebitur
\[
\begin{aligned}
A^{(h+1)} & = (-1)^h\,\frac{2}{\BareRoot{(1-nn)}}\left(\frac{1-\BareRoot{(1-nn)}}{n}\right)^{h-1}\\&\quad .  -\left(\frac{1-\BareRoot{(1-nn)}}{n}\right)^{2}\\&= (-1)^{h+1} . \frac{2}{\BareRoot{(1-nn)}}\left(\frac{1-\BareRoot{(1-nn)}}{n}\right)^{h+1},
\end{aligned}
\]
unde legis universalitas perspicitur.

